\chapter{Conclusion}

Travel OpSec isn't about fear, it's about intention. You
decide what to bring, what to leave behind, and how much
risk you're willing to carry. The less you expose, the
less you have to defend. Harden what matters, strip
everything else, and stay aware of your surroundings.

Security on the road isn't a checklist you tick once. It's
a posture you maintain. Updates, encryption, strong
authentication, minimal data, trusted paths -- these are
your baseline. From there, you layer additional controls
according to your context. Some trips need burner devices
and strict compartmentalization. Others just need good
hygiene and awareness. The point is to make those
decisions consciously, not out of habit or convenience.

When you return, assume that anything exposed has to be
re-evaluated. Rotate credentials, wipe or rebuild if
needed, and adjust your threat model based on what you
experienced. Each trip should make you sharper, not more
complacent.

Operational security -- and I'd dare to say security in
general -- isn't about perfection. It's about making it
harder for others to exploit you, and easier for you to
recover if something happens. Keep the footprint small,
act deliberately, and move through the world with clear
eyes.
